\section{量子位相推定}
\label{sec:quantum-phase-estimation}

\subsection{問題の設定}
\label{subsec:qpe-problem-setting}

ユニタリ演算子$U$の固有状態$\ket{u}$と固有値の関係を
\begin{equation}
  U\ket{u}
  =
  e^{2\pi\mathrm{i}\phi}\ket{u},
  \qquad
  0\leq\phi<1
  \label{eq:qpe-eigenvalue-equation}
\end{equation}
とする．
ユニタリ演算子の固有値は絶対値が$1$であるため，
固有値に含まれる情報は位相$\phi$として表すことができる．
\textbf{量子位相推定}は，$U$と$\ket{u}$を入力とし，
$\phi$の$t$ビット近似$\widetilde{\phi}$を補助量子ビットへ書き出すサブルーチンである
\cite{kitaev1995,shimada2020}．

量子位相推定では，$t$量子ビットの位相レジスタと，
固有状態$\ket{u}$を保持する状態レジスタを使用する．
初期状態は
\begin{align*}
  \ket{0}^{\otimes t}\ket{u}
\end{align*}
である．
位相レジスタの量子ビット数$t$を増やすほど，位相を細かく区別できる．

\subsection{位相キックバック}
\label{subsec:qpe-phase-kickback}

まず，一つの補助量子ビットを$\ket{+}$に準備し，
これを制御量子ビットとして制御$U$演算を適用する．
式\eqref{eq:qpe-eigenvalue-equation}より
\begin{align}
  \frac{1}{\sqrt{2}}
  (\ket{0}+\ket{1})\ket{u}
  &\longmapsto
  \frac{1}{\sqrt{2}}
  \left(
    \ket{0}\ket{u}
    +\ket{1}U\ket{u}
  \right)\notag\\
  &=
  \frac{1}{\sqrt{2}}
  \left(
    \ket{0}
    +e^{2\pi\mathrm{i}\phi}\ket{1}
  \right)\ket{u}
  \label{eq:single-qubit-phase-kickback}
\end{align}
となる．
固有状態$\ket{u}$自体は変化せず，固有値の位相だけが制御量子ビットの相対位相へ移る．
この現象を\textbf{位相キックバック}という．

同様に，制御$U^{2^r}$演算を適用すると
\begin{align*}
  \frac{1}{\sqrt{2}}(\ket{0}+\ket{1})\ket{u}
  \longmapsto
  \frac{1}{\sqrt{2}}
  \left(
    \ket{0}+e^{2\pi\mathrm{i}2^r\phi}\ket{1}
  \right)\ket{u}
\end{align*}
となる．
指数$2^r$を変えることで，$\phi$の2進表現に含まれる異なる桁の情報を取り出せる．

\subsection{制御ユニタリ演算による位相の書き込み}
\label{subsec:qpe-controlled-unitaries}

位相レジスタの各量子ビットへHadamardゲートを作用させると
\begin{align*}
  \ket{0}^{\otimes t}
  \longmapsto
  \frac{1}{\sqrt{M}}
  \sum_{k=0}^{M-1}\ket{k},
  \qquad
  M\coloneq2^t
\end{align*}
となる．
次に，位相レジスタの各量子ビットを制御部とする
$U,U^2,U^{2^2},\ldots,U^{2^{t-1}}$を状態レジスタへ作用させる．
計算基底状態$\ket{k}$が選択する演算は$U^k$であるため，レジスタ全体は
\begin{align}
  \frac{1}{\sqrt{M}}
  \sum_{k=0}^{M-1}\ket{k}\ket{u}
  &\longmapsto
  \frac{1}{\sqrt{M}}
  \sum_{k=0}^{M-1}\ket{k}U^k\ket{u}\notag\\
  &=
  \frac{1}{\sqrt{M}}
  \sum_{k=0}^{M-1}
  e^{2\pi\mathrm{i}k\phi}\ket{k}\ket{u}
  \label{eq:qpe-phase-register-state}
\end{align}
となる．
これにより，固有位相$\phi$が位相レジスタ全体の相対位相へ書き込まれる．

\subsection{逆量子フーリエ変換による読み出し}
\label{subsec:qpe-inverse-qft-readout}

まず，$\phi$が$t$ビットで厳密に表される場合を考える．
ある整数$m\in\{0,1,\ldots,M-1\}$を用いて
\begin{equation}
  \phi
  =
  \frac{m}{M}
  =
  0.\phi_1\phi_2\cdots\phi_t
  \label{eq:qpe-exact-binary-phase}
\end{equation}
と表せるとする．
このとき式\eqref{eq:qpe-phase-register-state}の位相レジスタは，
式\eqref{eq:qft-definition}の$F_M\ket{m}$と一致する．
したがって，逆量子フーリエ変換を作用させると
\begin{equation}
  \frac{1}{\sqrt{M}}
  \sum_{k=0}^{M-1}e^{2\pi\mathrm{i}km/M}\ket{k}\ket{u}
  \xmapsto{F_M^{\dagger}}
  \ket{m}\ket{u}
  \label{eq:qpe-exact-output}
\end{equation}
となる．
位相レジスタを計算基底で測定すれば，$\phi$の2進表現
$\phi_1\phi_2\cdots\phi_t$が確率$1$で得られる．

量子位相推定の処理は，次の手順にまとめられる．

\begin{enumerate}[(1)]
  \item 位相レジスタを$\ket{0}^{\otimes t}$，状態レジスタを$\ket{u}$に準備する．
  \item 位相レジスタへHadamardゲートを作用させ，一様な重ね合わせを生成する．
  \item 制御$U^{2^r}$演算によって，固有位相を位相レジスタへ書き込む．
  \item 位相レジスタへ逆量子フーリエ変換を作用させる．
  \item 位相レジスタを測定し，固有位相の近似値を得る．
\end{enumerate}

\subsection{有限ビット精度での推定}
\label{subsec:qpe-finite-precision}

一般の位相$\phi$は，有限の$t$ビットで厳密に表せるとは限らない．
式\eqref{eq:qpe-phase-register-state}へ逆量子フーリエ変換を作用させた後，
位相レジスタで整数$y$を得る確率振幅を
\begin{equation}
  \alpha_y
  \coloneq
  \frac{1}{M}
  \sum_{k=0}^{M-1}
  e^{2\pi\mathrm{i}k(\phi-y/M)}
  \label{eq:qpe-finite-amplitude}
\end{equation}
と定義する．
$\delta_y\coloneq\phi-y/M$とおくと，有限等比級数より
\begin{equation}
  |\alpha_y|^2
  =
  \frac{1}{M^2}
  \frac{\sin^2(\pi M\delta_y)}
       {\sin^2(\pi\delta_y)}
  \label{eq:qpe-probability-distribution}
\end{equation}
となる．
この確率分布は，$y/M$が$\phi$に近い値の周囲へ集中する．
したがって，測定値$y$から
\begin{equation}
  \widetilde{\phi}
  \coloneq
  \frac{y}{2^t}
  \label{eq:qpe-phase-estimate}
\end{equation}
を固有位相の近似値として得る．

位相レジスタを1量子ビット増やすと，表現できる位相の間隔は半分になる．
一方で，制御するユニタリ演算の最大指数は2倍になる．
したがって，位相精度を高めるには，補助量子ビット数だけでなく，
高いべき乗の制御$U$を実行する回路の深さと誤差も考慮する必要がある．

\subsection{固有状態の重ね合わせに対する作用}
\label{subsec:qpe-superposition-input}

実際の量子アルゴリズムでは，状態レジスタが一つの固有状態であるとは限らない．
$U$の正規直交固有状態を$\ket{u_j}$とし，一般の入力状態を
\begin{equation}
  \ket{\psi}
  \coloneq
  \sum_j c_j\ket{u_j}
  \label{eq:qpe-general-input}
\end{equation}
とする．
量子位相推定は線形であるため
\begin{equation}
  \ket{0}^{\otimes t}\ket{\psi}
  \longmapsto
  \sum_j c_j
  \ket{\widetilde{\phi}_j}\ket{u_j}
  \label{eq:qpe-superposition-output}
\end{equation}
と作用する．
位相レジスタを測定すると，結果$\widetilde{\phi}_j$はおよそ$|c_j|^2$の確率で得られ，
状態レジスタは対応する$\ket{u_j}$へ射影される．
一方，測定を行わずに位相レジスタを制御情報として用いれば，
すべての固有状態成分を重ね合わせのまま処理できる．
HHLアルゴリズムでは，この性質を用いて固有値ごとに異なる回転を作用させる．

\subsection{エルミート行列の固有値推定}
\label{subsec:qpe-hermitian-eigenvalues}

量子位相推定の入力$U$はユニタリ演算子でなければならない．
一方，線形方程式で扱う行列$A$は一般にユニタリではない．
$A$がエルミート行列であれば，実数$t_0$に対して
\begin{equation}
  U_A
  \coloneq
  e^{\mathrm{i}At_0}
  \label{eq:qpe-hamiltonian-unitary}
\end{equation}
はユニタリ演算子となる．
$A\ket{u_j}=\lambda_j\ket{u_j}$ならば
\begin{align*}
  U_A\ket{u_j}
  =
  e^{\mathrm{i}\lambda_jt_0}\ket{u_j}
  =
  e^{2\pi\mathrm{i}\phi_j}\ket{u_j}
\end{align*}
であるため，固有位相を
\begin{equation}
  \phi_j
  \coloneq
  \frac{\lambda_jt_0}{2\pi}
  \pmod{1}
  \label{eq:eigenvalue-to-phase}
\end{equation}
として推定できる．
固有値$\lambda_j$を一意に復元するには，固有値の既知の範囲に合わせて$t_0$を選び，
異なる固有値が同じ位相へ折り返されないようにする必要がある．

\subsection{逆量子位相推定}
\label{subsec:inverse-quantum-phase-estimation}

式\eqref{eq:qpe-superposition-output}の位相レジスタを用いて別の制御演算を実行した後，
量子位相推定回路を逆順に適用すると，位相レジスタを$\ket{0}^{\otimes t}$へ戻せる．
これを\textbf{逆量子位相推定}という．
位相の情報が目的のレジスタへ既に反映されていれば，
逆量子位相推定後もその効果は残り，位相レジスタだけが初期化される．
HHLアルゴリズムでは，固有値に応じた制御回転の後に逆量子位相推定を行い，
固有値レジスタと解状態との不要な量子もつれを取り除く．
